%=====================================================================
% Point-by-point response to the Review-03 comments (Group 09).
% Compile with pdflatex. Numbers quoted here are those of the revised
% manuscript (paper/main.tex, tables generated from outputs/review03).
%=====================================================================
\documentclass[11pt]{article}
\usepackage[margin=2.2cm]{geometry}
\usepackage{amsmath,amssymb}
\usepackage{booktabs}
\usepackage{enumitem}
\usepackage{xcolor}
\usepackage[hidelinks]{hyperref}
\newcommand{\Lspec}{L_{\mathrm{spec}}}
\newenvironment{comment}{\par\medskip\noindent\begin{minipage}{\linewidth}\itshape\color{black!70}}{\end{minipage}\par\smallskip}
\setlength{\parindent}{0pt}
\setlength{\parskip}{4pt}

\begin{document}

\begin{center}
{\Large\bfseries Response to Review-03 comments}\\[4pt]
Characteristic-Frequency-Guided Spectral Priors for Rolling-Element Bearing Fault Diagnosis\\
V.~Verma, T.~Garg, S.~Kumar, K.~P (Group 09)
\end{center}

We thank the reviewer for a careful and constructive assessment. Every
criticism was valid, and addressing them changed the paper's conclusions. In
summary:
\begin{itemize}[leftmargin=*]
\item All results now come from the \textbf{original CWRU recordings} (40 files)
and a \textbf{second dataset (MFPT)}. The constructed dataset is no longer used
for any result.
\item Every split is formed \textbf{before windowing}, and an automated check
confirms that no two windows in different splits share a sample.
\item The spectral-consistency term was \textbf{audited} and found to be defective
in three respects; a corrected version (v2) is derived and unit-tested.
\item The method is repositioned as a \textbf{physics-guided spectral prior},
explicitly distinguished from differential-equation PINNs.
\item All mandatory (P1) experiments were run with \textbf{five seeds}, 95\%
confidence intervals and paired tests. Of the suggested (P2) items, the
noise/speed, raw-vs-envelope and band-width/harmonic sensitivity studies are
complete; the sensitivity is measured on the prior's evidence accuracy without
retraining. The MFPT classifier results rest on 2--3 seeds and are labelled
preliminary, but the MFPT evidence analysis is complete and explains them;
the table in Section~3 states this explicitly.
\item The measured outcome is \textbf{negative}. On real data the spectral prior
does not improve cross-load accuracy over the identical network without it, and
its transductive variant degrades accuracy significantly. The title, abstract,
contributions and conclusions have been rewritten to say so.
\item The complete training and evaluation code, configurations and per-seed run
files are now part of the submission.
\end{itemize}

%---------------------------------------------------------------------
\section*{1. Critical methodological weaknesses}

\begin{comment}
1. The central research claim remains untested. [\dots] The original CWRU
recordings must be used to establish the proposed contribution.
\end{comment}
\textbf{Response.} Agreed. We downloaded the 40 standard 12\,kHz drive-end
recordings from the CWRU Bearing Data Center: normal baselines 97--100, and
inner-race, ball and outer-race (@6:00) faults of 0.007/0.014/0.021\,in at
0--3\,hp. The 48\,kHz normal baselines are decimated to 12\,kHz, and the recorded
shaft speed of each file is used for its characteristic frequencies. Every
diagnosis result in the revised paper (Tables~VI--XII) is measured on these
recordings. \emph{Changes:} Sections~IV and VII, Tables~III--XII.

\begin{comment}
2. The synthetic experiment risks circular validation. [\dots]
\end{comment}
\textbf{Response.} Agreed. The constructed signals were generated from the same
defect-frequency relations that the loss rewards, so success on them only showed
compatibility with the assumed model. We removed the constructed data from all
evaluations and added a paragraph (Section~IV-D) explaining why it cannot support
a robustness claim. The real-data results confirm the reviewer's concern: on real
recordings the prior conflicts with the labels for a large share of windows,
which never happened on the constructed data. One example: envelope evidence alone
identifies the fault type of only 66\% of the 0--2\,hp training windows (ball
faults 50\%). \emph{Changes:} Section~IV-D; Table~IV.

\begin{comment}
3. The spectral-consistency loss needs closer examination. [\dots] verify that
this normalization, the healthy-class zero mask and gradient propagation produce
the intended optimization behaviour.
\end{comment}
\textbf{Response.} We performed this verification. It showed that the original
term did \emph{not} behave as intended (Section~V-B, Table~IV):
\begin{enumerate}[leftmargin=*]
\item \emph{Normalised frequency input.} The term used the time-summed $64\times64$
STFT image, whose rows are 95\,Hz apart. All defect harmonics (up to 487\,Hz)
fall in its lowest six rows, and log compression flattens them. Used on its own,
it identifies the fault type of 32\% of fault windows, which is chance level. Its
gradient is about $10^{-6}$ per logit, so the adaptive weight inflated to
$\bar\lambda\approx40$ during training.
\item \emph{Healthy-class zero mask and gradient.} The zero mask does not make the
term neutral to the healthy class. Its window weight
$w_b=\sum_{c\in\text{faults}}p_{bc}$ is differentiable, so windows that align
worse than the batch average are pushed toward ``healthy''. This happens for
64\% of the faulty training windows.
\item \emph{Mixture normalisation.} Normalising the probability-weighted mask
mixture couples each class's pull to its discretised mask mass (3.2 for
inner/outer vs.\ 2.3 for ball).
\end{enumerate}
The revised term (v2, Eq.~(7)) makes three changes:
\begin{itemize}[leftmargin=*]
\item it uses a 5.9\,Hz-resolution raw or squared-envelope spectrum computed from the window;
\item it gives every class mask the same total mass;
\item it uses the fault-type softmax over fault logits only, with a stop-gradient window weight.
\end{itemize}
As a result $\partial\Lspec/\partial z_{\mathrm{healthy}}=0$ exactly, and among
fault types probability moves toward the type with stronger evidence. Both
properties are checked by unit tests (\texttt{tests/test\_review03.py}). The
original term is kept only as an ablation row.

\begin{comment}
4. The physics term may not enforce a genuinely independent physical constraint.
[\dots] the paper should distinguish it from conventional differential-equation PINNs.
\end{comment}
\textbf{Response.} Agreed, and stated explicitly in a new subsection
(Section~V-D). The term combines a kinematic relation $f_c=r_c f_r$ with the
observed spectrum of the same input. It is not the residual of a governing
equation and does not encode resonance dynamics, slip or modulation. We call it
a physics-\emph{guided spectral prior} throughout and removed ``PINN'' from the
title and the claims. The closest analogue of PINN collocation points is to apply
the prior to unlabelled target-condition windows. We evaluated this variant
separately and labelled it transductive. It reduced 3\,hp accuracy from
99.25\,$\pm$\,0.95\% to 89.63\,$\pm$\,0.49\% ($p<0.001$), because the envelope
evidence is wrong for many target windows. The constraint is therefore only as
reliable as the spectral evidence it relies on.

\begin{comment}
5. Generalization is not established. A zero accuracy drop on constructed data
cannot demonstrate robustness when all competing neural models achieve
essentially perfect accuracy.
\end{comment}
\textbf{Response.} Agreed. We no longer claim robustness from the constructed
data. On real data we report mean\,$\pm$\,95\% CI over five seeds and paired
$t$-tests against the identical network without physics. We also added harder
settings where models do not saturate: single-source 0$\to$3\,hp, additive noise
down to $-5$\,dB, emulated speed change, and MFPT. Under the main 0/1/2$\to$3\,hp
protocol the CNN--LSTM reaches 99.25\,$\pm$\,0.95\%, and no physics variant
improves on it significantly. In the robustness tests the largest gain from the
prior is $+3.4$ points at a $\times1.05$ speed change ($p=0.08$, not significant).
The paper now states plainly that generalisation benefits are not established.

%---------------------------------------------------------------------
\section*{2. Dataset and implementation audit}

\begin{comment}
No executable training script or notebook is included in this archive. [\dots]
The manuscript should not describe the implementation as fully reproducible
until the complete training and evaluation code is provided.
\end{comment}
\textbf{Response.} The revised archive contains the complete repository:
\begin{itemize}[leftmargin=*]
\item data loaders and protocols (\texttt{src/review03/data\_real.py});
\item losses (\texttt{src/review03/spectral.py});
\item training loop (\texttt{src/review03/train.py});
\item experiment runner (\texttt{scripts/run\_review03.py});
\item table and figure generators;
\item unit tests;
\item one JSON file per trained run.
\end{itemize}
Every number in the result tables is generated from these run files, with no
hand-typed values. Section~X gives the exact commands and the data sources,
including checksums. The paper now notes that GPU (MPS) kernels are not bit-wise
deterministic.

%---------------------------------------------------------------------
\section*{3. Mandatory work for Review-03}
\begin{center}
\small
\begin{tabular}{p{0.6cm}p{4.6cm}p{2.0cm}p{8.0cm}}
\toprule
Pri. & Required experiment & Status & Outcome (revised paper) \\
\midrule
P1 & Original CWRU recordings & Done & 40 recordings, 4 classes; Tables VI--X. \\
P1 & Train 0/1/2\,hp, test 3\,hp & Done & CNN--LSTM 99.3\,$\pm$\,1.0\%; no physics variant significantly better (Table VII). Harder 0$\to$3\,hp added: 87.1\,$\pm$\,6.2\%, again no significant gain (Table IX). \\
P1 & CNN--LSTM with vs.\ without physics loss & Done & Identical network, $\lambda=0$ vs.\ v1/v2 raw/v2 envelope/transductive; paired $t$-tests. Envelope $-1.6$ pts ($p=0.26$); transductive $-9.6$ pts ($p<0.001$). \\
P1 & Fixed vs.\ adaptive weighting & Done & Fixed $\lambda=0.1$ 99.4\%, 0.5 99.1\%, 1.0 98.4\%, adaptive 97.6\%; adaptive balancing is not supported (Table VIII). \\
P1 & Repeated runs, confidence intervals & Done & 5 seeds, mean\,$\pm$\,95\% CI, paired $t$-test vs.\ baseline in every CWRU table. \\
\midrule
P2 & Band-width / harmonic sensitivity & Done (training-free) & Evidence accuracy over $\sigma\in\{2,5,10,20\}$\,Hz, $H\in\{1,3,5\}$: envelope 51--67\% on CWRU 3\,hp, default within 3.1 pts of best; raw $\approx$ chance. The prior is insensitive to its settings but never reliable on CWRU (Table~XII, Sec.~VII-I). \\
P2 & Noisy and speed-varying signals & Done (speed emulated) & Noise 10 to $-5$\,dB and $\times0.95$/$\times1.05$/$\pm5\%$ time-warp; no significant gain from the prior; best case $+3.4$ pts at $\times1.05$, $p=0.08$ (Table X, Fig.~3). \\
P2 & Raw vs.\ envelope spectrum & Done & Physics-only fault-type accuracy 33\% (raw) vs.\ 66\% (envelope) (Table IV); in training, raw is neutral and envelope slightly harmful (Table VII). \\
P2 & Another bearing dataset & Done (classifier runs preliminary) & MFPT load shift: supervised models 29--39\% due to a load/class confound; transductive envelope prior 100\% (2 seeds). Envelope evidence is 100\% correct on MFPT vs.\ 64\% on CWRU, which explains the reversal (Tables XI--XII). \\
\midrule
-- & Avoid window-level leakage & Done & Splits made before windowing, guard gaps, automated overlap check (Section IV-C, Table III). \\
\bottomrule
\end{tabular}
\end{center}

\textbf{Window-level leakage.} Splits are formed on the continuous recordings
before windowing:
\begin{itemize}[leftmargin=*]
\item \emph{cross-load protocols:} the whole 3\,hp recordings are the test set;
\item \emph{in-distribution:} each recording is cut into 70/15/15\% contiguous
blocks with a one-window guard gap.
\end{itemize}
\texttt{assert\_no\_overlap} raises an error if any two windows in different
splits share a sample, and it runs on every protocol (Table~III). CWRU has one
recording per class, severity and load. A by-recording split within one load
would therefore remove a whole fault condition from training, so the
in-distribution protocol uses time blocks.

%---------------------------------------------------------------------
\section*{4. Summary of changes to the manuscript}
\begin{itemize}[leftmargin=*]
\item Title, abstract, contributions and conclusion rewritten to match the
measured (negative) results.
\item Removed: constructed-dataset results, literature-anchored ``target''
columns, and the claim of full reproducibility without code.
\item New sections:
  \begin{itemize}
  \item data and leakage-controlled protocols (IV);
  \item audit and revision of $\Lspec$ (V-B, V-C);
  \item physics-guided prior vs.\ PINN (V-D);
  \item limitations and threats to validity (VIII);
  \item reproducibility (X).
  \end{itemize}
\item New tables and figures:
  \begin{itemize}
  \item protocols, the $\Lspec$ audit, in-distribution, cross-load and single-source results;
  \item weighting, robustness, sensitivity and MFPT results;
  \item real-data envelope spectra.
  \end{itemize}
\item Added references: Zhao \emph{et al.} (2020) on evaluation leakage; MFPT data.
\end{itemize}

\end{document}
